% Part: first-order-logic
% Chapter: axiomatic-deduction
% Section: provability-propositional

\documentclass[../../../include/open-logic-section]{subfiles}

\begin{document}

\iftag{FOL}
      {\olfileid{fol}{axd}{ppr}}
      {\olfileid{pl}{axd}{ppr}}

\olsection{\usetoken{S}{derivability} lan Konektif Proposisional}

\begin{explain}
  Kita netepake yen relasi !!{derivability}~$\Proves$ saka deduksi
  aksiomatik cukup kuwat kanggo netepake sawetara kasunyatan dhasar
  kang ngemot konektif proposisional, kayata $!A \land !B \Proves !A$
  lan $!A, !A \lif !B \Proves !B$ (modus ponens). Kasunyatan kasebut
  dibutuhake ing bukti teorema kelengkapan.
\end{explain}

\begin{prop}\ollabel{prop:provability-land}
  \begin{enumerate}
  \item \ollabel{prop:provability-land-left} $!A \land !B \Proves !A$
    lan $!A \land !B \Proves !B$ kalorone lumaku.
  \item \ollabel{prop:provability-land-right} $!A, !B \Proves !A \land !B$.
  \end{enumerate}
\end{prop}

\begin{proof}
  \begin{enumerate}
    \item Saka \olref[prp]{ax:land1} lan \olref[prp]{ax:land2}
      nganggo modus ponens. Cathetan koreksi sumber OLPL-032: rujukan
      kapindho dibalekake marang aksioma proyeksi tengen, kang menehi
      kesimpulan B.
  \item Saka \olref[prp]{ax:land3} nganggo rong aplikasi modus ponens.
  \end{enumerate}
\end{proof}


\begin{prop}\ollabel{prop:provability-lor}
  \begin{enumerate}
  \item $!A \lor !B, \lnot !A, \lnot !B$ ora konsisten.
  \item $!A \Proves !A \lor !B$ lan $!B \Proves !A \lor !B$
    kalorone lumaku.
  \end{enumerate}
\end{prop}

\begin{proof}
  \begin{enumerate}
  \item Saka \olref[prp]{ax:lnot2}, kita entuk $\Proves \lnot !A
    \lif (!A \lif \lfalse)$ lan $\Proves \lnot !B \lif (!B \lif
    \lfalse)$. Cathetan koreksi sumber OLPL-033: wujud-wujud kasebut
    minangka conto aksioma negasi kapindho, dudu aksioma
    kontraposisi. Mula, miturut Teorema Deduksi, kita duwe $\{\lnot
    !A\} \Proves !A \lif \lfalse$ lan $\{\lnot !B\} \Proves !B \lif
    \lfalse$. Saka \olref[prp]{ax:lor3}, kita entuk $\{\lnot !A,
    \lnot !B\} \Proves (!A \lor !B) \lif \lfalse$. Miturut teorema
    Deduksi, $\{!A \lor !B, \lnot !A, \lnot !B\} \Proves \lfalse$.
  \item Saka \olref[prp]{ax:lor1} lan \olref[prp]{ax:lor2} nganggo
    modus ponens.
  \end{enumerate}
\end{proof}

\begin{prop}\ollabel{prop:provability-lif}
  \begin{enumerate}
  \item \ollabel{prop:provability-lif-left} $!A, !A \lif !B \Proves !B$.
  \item \ollabel{prop:provability-lif-right}
    $\lnot !A \Proves !A \lif !B$ lan $!B \Proves !A \lif !B$
    kalorone lumaku.
  \end{enumerate}
\end{prop}

\begin{proof}
  \begin{enumerate}
  \item Kita bisa !!{derive}:
    \begin{derivation}
      1. & $!A$ & \Hyp\\
      2. & $!A \lif !B$ & \Hyp\\
      3. & $!B$ & 1, 2, \MP
    \end{derivation}
    
  \item Miturut \olref[prp]{ax:lnot2} lan \olref[prp]{ax:lif1} sarta
    Teorema Deduksi, miturut urutane.
  \end{enumerate}
\end{proof}

\end{document}
